%%%%%%%%%%%%%%%%%%%%%%%%%%%%%%%%%%%%%%%%%%%%%%%%%%%%%%%%%%%%%%%%%
% GIÁO ÁN STEM: CRYPTOCODE - ĐIỆP VIÊN MẬT MÃ
% Môn: Toán 10/11 | Chủ đề: Số học đồng dư & Mã hóa
%%%%%%%%%%%%%%%%%%%%%%%%%%%%%%%%%%%%%%%%%%%%%%%%%%%%%%%%%%%%%%%%%
\documentclass[12pt,a4paper]{article}

\usepackage[utf8]{inputenc}
\usepackage[T1]{fontenc}
\usepackage[vietnamese]{babel}
\usepackage{graphicx}
\usepackage{amsmath,amssymb}
\usepackage{array, booktabs, tabularx}
\usepackage{tikz}
\usetikzlibrary{calc}
\usepackage{geometry}
\usepackage{fancyhdr}
\usepackage[svgnames]{xcolor}
\usepackage{tcolorbox}
\tcbuselibrary{skins, breakable}
\usepackage{hyperref}
\usepackage{enumitem}

\geometry{left=2cm,right=2cm,top=2.5cm,bottom=2cm}
\definecolor{cryptogreen}{RGB}{16, 185, 129}   % Emerald 500
\definecolor{cryptodark}{RGB}{6, 78, 59}       % Emerald 900
\definecolor{cryptobg}{RGB}{236, 253, 245}     % Emerald 50

\hypersetup{colorlinks=true, linkcolor=cryptodark, urlcolor=cryptogreen}

% ============= BOX DEFINITIONS =============
\newtcolorbox{sectionbox}[1]{
    enhanced, breakable, colback=white, colframe=cryptogreen, boxrule=1pt,
    fonttitle=\bfseries\sffamily\Large, coltitle=white, colbacktitle=cryptogreen,
    title={#1},
    shadow={2mm}{-1mm}{0mm}{black!20},
    code={\phantomsection\addcontentsline{toc}{section}{#1}}
}

\newtcolorbox{activitybox}[1]{
    enhanced, breakable, colback=white, colframe=cryptodark, boxrule=1pt, arc=3mm,
    title={\sffamily\bfseries\color{white} [HOẠT ĐỘNG] #1},
    colbacktitle=cryptodark
}

\newtcolorbox{mathbox}[1]{
    enhanced, breakable, colback=cryptobg, colframe=cryptodark, boxrule=1pt, arc=3mm,
    title={\sffamily\bfseries\color{white} [TOÁN] #1},
    colbacktitle=cryptodark
}

% ============= HEADERS =============
\pagestyle{fancy}
\fancyhf{}
\fancyhead[L]{\color{cryptodark}\textbf{DỰ ÁN STEM: CRYPTOCODE}}
\fancyhead[R]{\color{cryptodark}\textbf{TOÁN HỌC \& BẢO MẬT}}
\fancyfoot[C]{\thepage}

\begin{document}

% ========== TITLE PAGE ==========
\begin{titlepage}
    \begin{tikzpicture}[remember picture, overlay]
        \fill[cryptodark] (current page.north west) rectangle (current page.south east);
        \node[anchor=center, color=cryptogreen, font=\fontsize{60}{70}\sffamily\bfseries] at ($(current page.center)+(0,1)$) {CRYPTO};
        \node[anchor=center, color=white, font=\fontsize{60}{70}\sffamily\bfseries] at ($(current page.center)+(0,-1)$) {CODE};
        \node[anchor=center, color=white, font=\Large\sffamily] at ($(current page.center)+(0,-3)$) {Mật mã học cho Thế hệ Số};
    \end{tikzpicture}
    
    \vspace{12cm}
    \begin{center}
        \color{black}
        \begin{tcolorbox}[width=0.8\textwidth, colback=white, colframe=cryptogreen]
            \centering \large
            \begin{tabular}{rl}
                \textbf{Môn học:} & Toán (Số học đồng dư) \\
                \textbf{Chủ đề:} & Phép Modulo, Mã Caesar \\
                \textbf{Ứng dụng:} & Bảo mật thông tin, An ninh mạng \\
                \textbf{Thời lượng:} & 2 tiết (90 phút)
            \end{tabular}
        \end{tcolorbox}
    \end{center}
\end{titlepage}

\tableofcontents
\newpage

% ========== CONTENT ==========

\begin{sectionbox}{I. MỤC TIÊU BÀI HỌC}
\phantomsection\subsection*{1. Kiến thức}\addcontentsline{toc}{subsection}{1. Kiến thức}
\begin{itemize}
    \item Hiểu khái niệm \textbf{Phép chia lấy dư} (Modulo) và ký hiệu $a \mod n$.
    \item Nắm được nguyên lý hoạt động của \textbf{Mã dịch chuyển Caesar}.
    \item Biết cách biểu diễn chữ cái thành số ($A=0, B=1, ..., Z=25$).
\end{itemize}

\phantomsection\subsection*{2. Kỹ năng}\addcontentsline{toc}{subsection}{2. Kỹ năng}
\begin{itemize}
    \item Thực hiện mã hóa và giải mã tin nhắn bằng công thức toán học.
    \item Sử dụng phân tích tần suất để phá mã.
    \item Lập trình đơn giản (nếu có điều kiện).
\end{itemize}

\phantomsection\subsection*{3. Phẩm chất}\addcontentsline{toc}{subsection}{3. Phẩm chất}
\begin{itemize}
    \item Ý thức bảo vệ thông tin cá nhân trong môi trường số.
    \item Tư duy logic, kiên nhẫn khi giải quyết vấn đề.
\end{itemize}
\end{sectionbox}

\begin{sectionbox}{II. CHUẨN BỊ}
\begin{itemize}
    \item \textbf{Giáo viên:} Máy chiếu, Web App \texttt{crypto.html}, Bảng chữ cái A-Z.
    \item \textbf{Học sinh:} Điện thoại/Laptop, Phiếu học tập, Bút chì.
\end{itemize}
\end{sectionbox}

\begin{sectionbox}{III. TIẾN TRÌNH DẠY HỌC (MÔ HÌNH 5E)}

\phantomsection\subsection*{1. ENGAGE (Gắn kết - 10p): Bức thư bí ẩn}\addcontentsline{toc}{subsection}{1. ENGAGE (Gắn kết - 10p)}
\textbf{Hoạt động:} GV đưa ra bức thư mật:
\begin{center}
\texttt{WKLV LV D VHFUHW PHVVDJH}
\end{center}
\textbf{Câu hỏi:} "Đây là một tin nhắn bí mật từ 2000 năm trước. Các em có thể giải mã không?"

\textbf{Dẫn dắt:} Julius Caesar đã dùng phương pháp này để liên lạc quân sự. Chỉ cần một con số bí mật (khóa), ta có thể mã hóa mọi thông điệp.

\phantomsection\subsection*{2. EXPLORE (Khám phá - 20p): Tự tay mã hóa}\addcontentsline{toc}{subsection}{2. EXPLORE (Khám phá - 20p)}
\begin{activitybox}{Bảng dịch chuyển Caesar (Khóa = 3)}
\begin{center}
\begin{tabular}{|c|c|c|c|c|c|c|c|c|c|c|c|c|}
\hline
A & B & C & D & E & F & G & H & I & J & K & L & M \\
\hline
D & E & F & G & H & I & J & K & L & M & N & O & P \\
\hline
\end{tabular}
\vspace{0.3cm}

\begin{tabular}{|c|c|c|c|c|c|c|c|c|c|c|c|c|}
\hline
N & O & P & Q & R & S & T & U & V & W & X & Y & Z \\
\hline
Q & R & S & T & U & V & W & X & Y & Z & A & B & C \\
\hline
\end{tabular}
\end{center}
HS thực hành:
\begin{enumerate}
    \item Mã hóa tên mình với khóa $k=3$.
    \item Giải mã: \texttt{WKLV LV D VHFUHW PHVVDJH} $\to$ ?
\end{enumerate}
\end{activitybox}

\phantomsection\subsection*{3. EXPLAIN (Giải thích - 25p): Toán học đằng sau}\addcontentsline{toc}{subsection}{3. EXPLAIN (Giải thích - 25p)}
\begin{mathbox}{Công thức Mã hóa Caesar}
Gán mỗi chữ cái một số: $A=0, B=1, ..., Z=25$.

\textbf{Mã hóa:} $E(x) = (x + k) \mod 26$

\textbf{Giải mã:} $D(y) = (y - k) \mod 26$

Trong đó:
\begin{itemize}
    \item $x$: Giá trị số của chữ cái gốc.
    \item $k$: Khóa dịch chuyển (số bước).
    \item $\mod 26$: Lấy dư khi chia cho 26 (quay vòng bảng chữ cái).
\end{itemize}

\textbf{Ví dụ:} Mã hóa chữ "H" (vị trí 7) với $k=3$:
$$ E(7) = (7 + 3) \mod 26 = 10 \to \text{K} $$
\end{mathbox}

\phantomsection\subsection*{4. ELABORATE (Áp dụng - 25p): Phá mã như Điệp viên}\addcontentsline{toc}{subsection}{4. ELABORATE (Áp dụng - 25p)}
\begin{activitybox}{Brute-force Attack}
Nếu không biết khóa $k$, ta có thể thử tất cả 26 giá trị có thể (từ 0 đến 25).

HS sử dụng App \textbf{CryptoCode} tại \href{https://stem-1h8.pages.dev/crypto.html}{\texttt{stem-1h8.pages.dev/crypto.html}}:
\begin{enumerate}
    \item Nhập tin mã hóa.
    \item Nhấn "Brute-force" để xem tất cả 26 phương án giải mã.
    \item Tìm câu có nghĩa.
\end{enumerate}

\textbf{Nâng cao:} Phân tích tần suất chữ cái (chữ E xuất hiện nhiều nhất trong tiếng Anh).
\end{activitybox}

\phantomsection\subsection*{5. EVALUATE (Đánh giá - 10p)}\addcontentsline{toc}{subsection}{5. EVALUATE (Đánh giá - 10p)}
HS hoàn thành thử thách:
\begin{enumerate}
    \item Tạo tin nhắn bí mật cho bạn cùng lớp với khóa tự chọn.
    \item Giải mã tin nhắn nhận được từ bạn khác.
\end{enumerate}
\end{sectionbox}

\newpage
\begin{sectionbox}{IV. PHỤ LỤC}
\phantomsection\subsection*{Phụ lục 1: Lịch sử Mật mã}\addcontentsline{toc}{subsection}{Phụ lục 1: Lịch sử Mật mã}
\begin{itemize}
    \item \textbf{Caesar Cipher (50 TCN):} Dịch chuyển bảng chữ cái.
    \item \textbf{Vigenère Cipher (TK 16):} Dùng từ khóa thay vì số đơn lẻ.
    \item \textbf{Enigma (WW2):} Máy mã hóa cơ điện phức tạp.
    \item \textbf{RSA (1977):} Mã hóa khóa công khai hiện đại (dùng số nguyên tố lớn).
\end{itemize}

\phantomsection\subsection*{Phụ lục 2: Ứng dụng thực tế}\addcontentsline{toc}{subsection}{Phụ lục 2: Ứng dụng thực tế}
\begin{itemize}
    \item HTTPS (Bảo mật website).
    \item Mã hóa đầu cuối (End-to-End Encryption) trong Zalo, Messenger.
    \item Chữ ký số, Blockchain.
\end{itemize}
\end{sectionbox}

\end{document}
